% !TEX program = pdflatex
% Mini-évaluation de reprise - Seconde -- SUJET B (outils de calcul, fractions, puissances)
% Compilation : pdflatex eval-reprise-calcul-seconde-B.tex
\documentclass[11pt,a4paper]{article}

\usepackage[utf8]{inputenc}
\usepackage[T1]{fontenc}
\usepackage[french]{babel}
\usepackage{lmodern}
\usepackage{amsmath,amssymb,amsfonts}
\usepackage[a4paper,margin=2.2cm]{geometry}
\usepackage{enumitem}
\usepackage{fancyhdr}
\usepackage{xcolor}
\usepackage{array}

%------------------------------------------------------------
% Mise en forme
%------------------------------------------------------------
\definecolor{bleu}{RGB}{30,80,140}

\pagestyle{fancy}
\fancyhf{}
\lhead{\small Mathématiques}
\rhead{\small Seconde \textbullet\ Sujet B}
\cfoot{\small \thepage}
\renewcommand{\headrulewidth}{0.4pt}

% Environnement "exercice" numéroté avec titre (+ barème à droite)
\newcounter{exo}
\newcommand{\exo}[1]{%
  \refstepcounter{exo}%
  \vspace{0.7cm}
  \noindent{\color{bleu}\large\bfseries Exercice \theexo\ -- #1}\par
  \nopagebreak\smallskip
}

\setlength{\parindent}{0pt}
\setlist[enumerate]{leftmargin=*,itemsep=4pt,topsep=4pt}

%------------------------------------------------------------
\begin{document}

%------------------- En-tête -------------------
\begin{center}
  {\color{bleu}\LARGE\bfseries Évaluation de reprise}\\[4pt]
  {\large Seconde}\\[6pt]
  {\itshape Outils de calcul, fractions et puissances}
\end{center}

\vspace{6pt}
\noindent\fbox{\parbox{\dimexpr\textwidth-2\fboxsep-2\fboxrule}{%
  \small \textbf{Durée : 20 minutes.}

  \textbf{Calculatrice interdite.} \quad

  \vspace{5pt}
  Détaillez vos calculs \underline{\textbf{sur une copie séparée}} et donnez chaque résultat sous sa forme la plus simple.}}

%================= FRACTIONS =================

\exo{Fractions\hfill\normalfont\normalsize (4 points)}
Donner chaque résultat sous forme de fraction \textbf{irréductible}.
\begin{enumerate}
  \item $A = \dfrac{5}{6} - \dfrac{3}{4}$ \qquad ; \qquad $B = \dfrac{4}{5} \times \dfrac{15}{8}$
  \qquad ; \qquad $C = \dfrac{7}{10} \div \dfrac{14}{5}$
  \item Simplifier : $D = \dfrac{45}{18}$
\end{enumerate}

%================= PUISSANCES =================

\exo{Puissances\hfill\normalfont\normalsize (6 points)}
\begin{enumerate}
  \item Écrire sous la forme d'une seule puissance :

  \vspace{5pt}
  \hspace{1cm} $A = 3^4 \times 3^2$ \hspace{1.2cm} $B = (5^3)^2$ \hspace{1.2cm} $C = \dfrac{2^9}{2^6}$
  \item Écrire sous forme d'une fraction (sans exposant négatif) : $D = 5^{-2}$
  \item Écrire en notation scientifique : $27\,800$ \quad et \quad $0{,}00054$
\end{enumerate}

%================= OUTILS DE CALCUL =================

\exo{Outils de calcul\hfill\normalfont\normalsize (4 points)}
\begin{enumerate}
  \item Développer et réduire :

  \vspace{5pt}
  \hspace{1cm} $A = (x - 2)(x + 6)$ \hspace{1.5cm} $B = (3x + 1)^2$
  \item Factoriser :

  \vspace{5pt}
  \hspace{1cm} $A = 25x^2 - 4$ \hspace{1.5cm} $B = 7x^2 + 14x$
\end{enumerate}

\end{document}
